%% Chapter 8 — Section 8.1: Hamiltonian Dynamics
%% Source: PDF pages 100–102 (book pages 87–89 + 90–91)

\section{Hamiltonian Dynamics}
\label{sec:8.1}

This chapter is organized as follows. Hamiltonian dynamics are reviewed in Section~\ref{sec:8.1},
emphasizing three properties that are important to the design of sampling algorithms. The HMC
algorithm and an idealized version of it are introduced in Section~\ref{sec:8.2}. We focus on
showing that these algorithms are exact, meaning that they leave the target distribution
invariant. Section~\ref{sec:8.3} contains bibliographical remarks, including references that
address the geometric ergodicity of HMC.

\medskip
Given a Hamiltonian $H(q,p)$, the corresponding Hamilton equations are
\begin{equation}
\label{eq:8.2}
\begin{aligned}
\frac{dq_i}{dt} &= \frac{\partial H}{\partial p_i}, & 1 \le i \le d, \\[4pt]
\frac{dp_i}{dt} &= -\frac{\partial H}{\partial q_i}, & 1 \le i \le d.
\end{aligned}
\end{equation}

We vectorize equations~\eqref{eq:8.2} as follows
\[
\frac{dx}{dt} = J^{-1}\nabla H(x), \qquad
J := \begin{bmatrix} 0 & -I_{d\times d} \\ I_{d\times d} & 0 \end{bmatrix}
\in \R^{2d\times 2d}, \quad
x := \begin{bmatrix} q \\ p \end{bmatrix} \in \R^{2d},
\]
where
\[
\nabla H = \left[\frac{\partial H}{\partial q_1}, \ldots,
\frac{\partial H}{\partial q_d},
\frac{\partial H}{\partial p_1}, \ldots,
\frac{\partial H}{\partial p_d}\right]^T, \qquad
\frac{d}{dt}\begin{bmatrix}q\\p\end{bmatrix}
= \left[\frac{dq_1}{dt},\ldots,\frac{dq_d}{dt},\frac{dp_1}{dt},\ldots,\frac{dp_d}{dt}\right]^T.
\]

The matrix $J$ plays a central role in the Hamiltonian formalism. Note that $J$ is
skew-symmetric, meaning that $J^T = -J$. Later we will use the following properties of
invertible skew-symmetric matrices:
\begin{enumerate}
  \item $J^{-1}$ is also skew-symmetric, since $J^{-1} = -(J^T)^{-1} = -(J^{-1})^T$.
  \item The corresponding quadratic form is zero, since $x^T J x = x^T J^T x = -x^T J x$.
\end{enumerate}

Hamilton equations~\eqref{eq:8.2} define a flow in phase space: For $t \ge 0$, we denote by
$\varphi_t : \R^{2d} \to \R^{2d}$ the $t$-flow map, so that $\varphi_t(x) \in \R^{2d}$ is
the solution to Hamilton equations at time $t$ with initial condition $x \in \R^{2d}$. It
follows directly from the definition that $\varphi_0(x) = x$, i.e.\ $\varphi_0$ is the
identity map. Furthermore, for $t, s \ge 0$ we have that $\varphi_{s+t}(x) = \varphi_s \circ
\varphi_t(x)$, since the solution at time $t + s$ with initial condition $x$ can be found by
first solving up to time $t$ to obtain $x(t) := \varphi_t(x)$, and then solving Hamilton
equations for an additional time $s$ with initial condition $x(t)$ to obtain
$x(t+s) := \varphi_s(x(t)) = \varphi_s \circ \varphi_t(x)$. The concept of flow map of a
dynamical system is helpful in exploring the qualitative behavior of numerical methods to solve
them, since one can think of numerical methods as approximations of the flow map. In particular,
we will see that in order to ensure that Hamiltonian Monte Carlo algorithms leave the extended
target invariant, it is essential to use numerical solvers that preserve key properties of the
flow map.

\begin{example}[Harmonic Oscillator]
Consider the case $d = 1$, $V(q) = \frac{q^2}{2}$, and $K(p) = \frac{p^2}{2}$. Then, the
Hamilton equations are
\[
\frac{dq}{dt} = p, \qquad \frac{dp}{dt} = -q.
\]
The general solutions have the form
\[
q(t) = r\cos(a + t), \qquad p(t) = -r\sin(a + t),
\]
for some constants $a$ and $r$, which shows that $\varphi_t$ is a clockwise rotation in the
$q-p$ plane.
\end{example}

Hamiltonian dynamics dovetail nicely with MCMC because they satisfy three important properties:
conservation of the Hamiltonian, symplecticity, and reversibility. We next explain these
properties and how they are relevant to the construction of MCMC algorithms.

%% -------------------------------------------------------
\subsection{Conservation of Hamiltonian}
\label{sec:8.1.1}
%% -------------------------------------------------------

By the chain rule and the fact that $J$ is skew-symmetric, we have that
\[
\frac{dH}{dt} = \nabla H^T \frac{dx}{dt} = \nabla H^T J^{-1} \nabla H = 0, \qquad
x = \begin{bmatrix}q\\p\end{bmatrix}.
\]

Recall the definition of $f^H$ in \eqref{eq:8.1}. The value of $f^H$ depends on $x = (q,p)$
only through $H$. Therefore, conservation of the Hamiltonian implies that if we start at any
$x_0 = (q_0, p_0)$ and move with Hamiltonian dynamics, we move along level sets of $f^H$. In
other words, we have $H(x_0) = H(\varphi_s(x_0))$ for all $s > 0$. This property, together
with the preservation of the canonical measure that we describe next, will allow us to use
Hamiltonian dynamics to propose bold moves that are still likely to be accepted.

\begin{figure}[htbp]
\centering
\includegraphics[width=0.60\textwidth]{figures/fig_08_1}
\caption{\textit{Movement along a level set of $f^H$ starting from $(q_0, p_0)$.}}
\label{fig:8.1}
\end{figure}

%% -------------------------------------------------------
\subsection{Symplecticity and Preservation of the Canonical Measure}
\label{sec:8.1.2}
%% -------------------------------------------------------

The flow map $x(s) \xrightarrow{\varphi_t} x(t+s)$ preserves volume. This is a direct
consequence of the symplecticity of Hamiltonian dynamics. We next recall the definition of
symplecticity and prove the symplecticity of Hamiltonian dynamics using a theorem by Poincar\'{e}.

\begin{definition}[Symplectic Maps]
A linear map $L : \R^{2d} \to \R^{2d}$ is symplectic if $L^T J L = J$. A continuously
differentiable map $\psi : \R^{2d} \to \R^{2d}$ is symplectic if its Jacobian is symplectic
everywhere. That is, if
\[
\psi'(x)^T J \psi'(x) = J, \qquad \forall x \in \R^{2d}.
\]
\end{definition}

\begin{remark}
In $\R^2$ symplecticity of linear maps is equivalent to area preservation of parallelograms;
in $\R^{2d}$ it is equivalent to conservation of the oriented areas of the projections of
parallelograms into the coordinate planes $(q_i, p_i)$.
\end{remark}

\begin{theorem}[Hamiltonian Flows are Symplectic]
Let $H$ be a $\mathcal{C}^2$ Hamiltonian. Then the corresponding flow map $\varphi_t$ is
symplectic for all $t$.
\end{theorem}

\begin{proof}
Let $\varphi(x,t) := \varphi_t(x)$ denote the flow map of time $t$ starting from $x$. We
need to show that, for any $t$ and arbitrary $x \in \R^{2d}$,
\begin{equation}
\label{eq:8.3}
\frac{\partial}{\partial x}\varphi(x,t)^T J \frac{\partial}{\partial x}\varphi(x,t) = J.
\end{equation}

First, the equality holds for $t = 0$, since $\varphi(x,0)$ is the identity map on $\R^{2d}$.
Therefore, the proof will be complete if we show that the left-hand side in
equation~\eqref{eq:8.3} is constant in time. To that end, recall that
$\frac{\partial}{\partial t}\varphi(x,t) = J^{-1}\nabla H(\varphi(x,t))$ and so, using that
$H$ is twice continuously differentiable, we have
\begin{align*}
\frac{\partial^2}{\partial t \partial x}\varphi(x,t)
&= \frac{\partial^2}{\partial x \partial t}\varphi(x,t) \\
&= \frac{\partial}{\partial x} J^{-1}\nabla H(\varphi(x,t)) \\
&= J^{-1}\nabla^2 H\!\left(\varphi(x,t)\right)\frac{\partial}{\partial x}\varphi(x,t),
\end{align*}
where $\nabla^2$ denotes the Hessian. This partial result and the skew-symmetry of $J$ give that
\begin{align*}
&\frac{\partial}{\partial t}\!\left(\frac{\partial}{\partial x}\varphi(x,t)^T J
  \frac{\partial}{\partial x}\varphi(x,t)\right) \\
&= \frac{\partial^2}{\partial t\partial x}\varphi(x,t)^T J
  \frac{\partial}{\partial x}\varphi(x,t)
  + \frac{\partial}{\partial x}\varphi(x,t)^T J
  \frac{\partial}{\partial t\partial x}\varphi(x,t) \\
&= \frac{\partial}{\partial x}\varphi(x,t)^T \nabla^2 H\!\left(\varphi(x,t)\right)
  J^{-T} J \frac{\partial}{\partial x}\varphi(x,t)
  + \frac{\partial}{\partial x}\varphi(x,t)^T J J J^{-1}\nabla^2 H\!\left(\varphi(x,t)\right)
  \frac{\partial}{\partial x}\varphi(x,t) \\
&= -\frac{\partial}{\partial x}\varphi(x,t)^T \nabla^2 H\!\left(\varphi(x,t)\right)
  \frac{\partial}{\partial x}\varphi(x,t)
  + \frac{\partial}{\partial x}\varphi(x,t)^T \nabla^2 H\!\left(\varphi(x,t)\right)
  \frac{\partial}{\partial x}\varphi(x,t) \\
&= 0,
\end{align*}
and the proof is complete.
\end{proof}

Taking the determinant at both sides of the equality
\[
\left(\varphi_t'\right)^T J \varphi_t = J
\]
shows that the Jacobian determinant of $\varphi_t$ has absolute value 1. Using this property
and the conservation of the Hamiltonian, we obtain the following important result:

\begin{theorem}[Preservation of Canonical Measure]
The flow $\varphi_t$ preserves the canonical measure of $H$.
\end{theorem}

\begin{proof}
Let $D$ be a Lebesgue measurable set. We have
\begin{align*}
\mu^H\!\left(\varphi_t^{-1}(D)\right)
&= \int_{\varphi_{-t}(D)} \frac{1}{Z} e^{-H(x)}\,dx \\
&= \int_D \frac{1}{Z} e^{-H(\varphi_{-t}(x))} \left|\varphi_{-t}'(x)\right| dx \\
&= \int_D \frac{1}{Z} e^{-H(x)}\,dx \\
&= \mu^H(D).
\qedhere
\end{align*}
\end{proof}

%% -------------------------------------------------------
\subsection{Time Reversibility}
\label{sec:8.1.3}
%% -------------------------------------------------------

The term time reversible comes from classical mechanics. Let $q(t)$ and $p(t)$ be the height
and momentum of a pendulum of mass $m$ at time $t$. Suppose we release the pendulum at
$q(0) = q_0$ with momentum $p(0) = v_0 m$ in a frictionless space. By conservation of energy,
we know that the pendulum will come back to the same position $q_0$ at some time $t$, and its
velocity will be $-v_0$. The $q-p$ time plot is symmetric i.e.\ the system is time reversible.
Here we give a definition of time reversibility in the context of dynamical systems.

\begin{definition}[Involution and Reversibility]
A linear map $S$ is an involution if $S^2 = S \circ S$ is the identity map. A bijection
$\varphi$ is called reversible with respect to an involution $S$ if $S \circ \varphi =
\varphi^{-1} \circ S$. A differential equation is called reversible with respect to an
involution $S$ if, for every fixed $t$, its flow $\varphi_t$ is reversible with respect to $S$.
\end{definition}

The inverse of the flow $\varphi_t$ of a differential equation $\frac{dx}{dt} = u(x)$ can be
intuitively viewed as the flow map $\varphi_t$ of the system $\frac{dx}{dt} = -u(x)$. We can
give an equivalent characterization of time reversibility:

\begin{theorem}[Characterization of Reversibility]
A system of differential equations $\frac{dx}{dt} = u(x)$ is time reversible with respect to
an involution $S$ if and only if $S \circ u = -u \circ S$.
\end{theorem}

\begin{proof}
Suppose $\frac{dx}{dt} = u(x)$ is time reversible with respect to $S$. Using linearity of
$S$, we have
\begin{align*}
S \circ u(x)
&= \lim_{\epsilon \to 0} \frac{1}{\epsilon}
  \Bigl(S \circ \varphi_{t+\epsilon}(x) - S \circ \varphi_t(x)\Bigr) \\
&= \lim_{\epsilon \to 0} \frac{1}{\epsilon}
  \Bigl(\varphi_{t+\epsilon}^{-1} \circ S(x) - \varphi_t^{-1} \circ S(x)\Bigr) \\
&= \frac{d}{dt}\varphi_t^{-1}\!\left(S(x)\right) \\
&= (-u) \circ S(x).
\end{align*}
The other direction follows directly from the remark above.
\end{proof}

\begin{corollary}[Reversibility of Hamiltonian Systems]
Hamiltonian systems with Hamiltonian $H(q,p) = V(q) + K(p)$ such that $K(p)$ is an even
function are reversible with respect to the momentum flip involution
$S : (q,p) \mapsto (q,-p)$.
\end{corollary}

\begin{example}[Gaussian Momentum Variables and Reversibility]
The kinetic energy $K(p) = \frac{1}{2}p^T M^{-1} p$ is an even function, since
\[
K(-p) = \frac{1}{2}(-p)^T M^{-1}(-p) = \frac{1}{2}p^T M^{-1} p = K(p).
\]
Therefore, Corollary~8.8 implies that Hamiltonian systems with Hamiltonian
$H(q,p) = V(q) + \frac{1}{2}p^T M^{-1}p$ are always time reversible with respect to
momentum flip. Notice that with the choice of kinetic energy $K(p) = \frac{1}{2}p^T M^{-1}p$,
the marginal of $f^H(q,p) \propto \exp(-H(q,p))$ over the momentum variables is
Gaussian $\mathcal{N}(0, M)$.
\end{example}
